\documentclass[11pt,a4paper]{article}

\usepackage{amsmath,amssymb,amsthm}
\usepackage{mathtools}
\usepackage{bbm}
\usepackage{booktabs}
\usepackage{microtype}
\usepackage[margin=2.5cm]{geometry}
\usepackage[colorlinks=true,linkcolor=blue,citecolor=blue]{hyperref}

\newtheorem{remark}{Remark}
\newcommand{\Q}{\mathbb{Q}}
\newcommand{\E}{\mathbb{E}}
\newcommand{\R}{\mathbb{R}}
\newcommand{\F}{\mathcal{F}}
\newcommand{\dd}{\,\mathrm{d}}
\newcommand{\dW}{\dd W^{\Q}}
\newcommand{\norm}[1]{\left\lVert#1\right\rVert}

\title{\textbf{Deep BSDE Pricing of Interest Rate Caplets\\
under the Three-Factor Cheyette Model}\\[0.5em]
\large A White Paper Chapter}
\author{}
\date{}

\begin{document}
\maketitle
\tableofcontents
\bigskip

% ─────────────────────────────────────────────────────────────────────────────
\section{Introduction}

We present a deep BSDE solver for European caplet pricing under the
three-factor exponential Cheyette model of Beyna, Chiarella \& Kang
\cite{BCK2012}. The method follows the framework of Han, Jentzen \& E
\cite{HJE2018} and the xVA extension of Gnoatto, Picarelli \& Reisinger
\cite{GPR2022}, specialised to a non-recursive, linearly-driven backward
SDE whose solution coincides with the caplet price. The contribution is
twofold: (i) a careful derivation of the caplet FBSDE and its discretisation
within the Cheyette state-space; (ii) an architecture designed to meet a
stringent accuracy threshold of $0.1$ basis points vega-adjusted error, which
is approximately two orders of magnitude tighter than standard benchmarks in
the literature.

Although the Cheyette caplet admits a closed-form Black-type formula
\cite{BCK2012}, the deep BSDE framework is motivated by three considerations:
it generalises immediately to instruments without analytical solutions
(Bermudan swaptions, callable exotics, xVA integrals); it produces the
$\Delta$-hedge $\widehat{Z}_t$ as a direct output; and it provides a
pathwise simulation engine for exposure calculations without nested Monte
Carlo.

% ─────────────────────────────────────────────────────────────────────────────
\section{Model: Three-Factor Exponential Cheyette}
\label{sec:model}

Under the risk-neutral measure $\Q$, the short rate is
\begin{equation}
  r(t) = f(0,t) + X_1(t) + X_2(t) + X_3(t) + X_4(t),
\end{equation}
where $f(0,t)$ is the initial instantaneous forward rate (taken as a flat
curve $f_0$ throughout) and the state vector
$X = (X_1, X_2, X_3, X_4)^\top \in \R^4$ satisfies
\begin{align}
  \dd X_1 &= \bigl(V_{11}^{(1)}(t) + V_{12}^{(1)}(t)\bigr)\dd t
             + c\,\dd W^{(1)}_t, \label{eq:X1}\\
  \dd X_2 &= \bigl(V_{21}^{(1)}(t) + V_{22}^{(1)}(t) - \lambda_1 X_2\bigr)\dd t
             + P_1(t)\,\dd W^{(1)}_t, \label{eq:X2}\\
  \dd X_3 &= \bigl(V_{11}^{(2)}(t) - \lambda_2 X_3\bigr)\dd t
             + P_2(t)\,\dd W^{(2)}_t, \label{eq:X3}\\
  \dd X_4 &= \bigl(V_{11}^{(3)}(t) - \lambda_3 X_4\bigr)\dd t
             + P_3(t)\,\dd W^{(3)}_t, \label{eq:X4}
\end{align}
with $X_0 = 0$. Here $W^{(1)}, W^{(2)}, W^{(3)}$ are independent standard
Brownian motions under $\Q$; $P_k(t) = a_1^{(k)} t + a_0^{(k)}$ are
linear-in-time volatility functions; $\lambda_k > 0$ are mean-reversion
speeds; and $c > 0$ is a constant volatility shift for the first factor. The
deterministic functions $V_{ij}^{(k)}(t)$ are integrated-variance accumulators
that enforce the $\Q$-drift required by the HJM no-arbitrage condition.

The state dynamics can be written compactly as
\begin{equation}
  \dd X_t = \mu(t, X_t)\dd t + \sigma(t)\dW_t, \qquad X_0 = 0,
  \label{eq:X}
\end{equation}
where $\mu: [0,T] \times \R^4 \to \R^4$ collects the drifts in
\eqref{eq:X1}--\eqref{eq:X4} and $\sigma(t) \in \R^{4 \times 3}$ is the
diffusion matrix
\begin{equation}
  \sigma(t) =
  \begin{pmatrix}
    c    & 0      & 0      \\
    P_1(t) & 0    & 0      \\
    0    & P_2(t) & 0      \\
    0    & 0      & P_3(t)
  \end{pmatrix}.
\end{equation}
The rank of $\sigma$ is three; the $(X_1, X_2)$ block shares Brownian
$W^{(1)}$, producing the off-diagonal covariance that drives the
mixed-derivative term in the pricing PDE.

Zero-coupon bond prices are exponential-affine in $X$:
\begin{equation}
  B(t, T; X_t) = \frac{B(0,T)}{B(0,t)}
    \exp\!\Bigl[
      -G_1^{(1)}(t,T) X_1 - G_2^{(1)}(t,T) X_2
      - G_1^{(2)}(t,T) X_3 - G_1^{(3)}(t,T) X_4
      - H(t,T)
    \Bigr],
  \label{eq:bond}
\end{equation}
where the $G$ functions are integrals of the volatility structure and
$H(t,T) > 0$ is a convexity adjustment. Explicit forms are given in
\cite{BCK2012}.

% ─────────────────────────────────────────────────────────────────────────────
\section{Caplet FBSDE Formulation}
\label{sec:fbsde}

\subsection{Caplet payoff and the forward measure connection}

A caplet fixing at $T_C$ and paying at $T_B = T_C + \Delta$ has payoff (at
$T_B$)
\begin{equation}
  \Delta \cdot \max\bigl(R(T_C, T_B) - K,\, 0\bigr),
\end{equation}
where $R(T_C, T_B) = \Delta^{-1}(B(T_C, T_B)^{-1} - 1)$ is the spot LIBOR
at $T_C$. Discounting to $T_C$ produces the terminal condition
\begin{equation}
  \phi(X_{T_C}) = B(T_C, T_B; X_{T_C}) \cdot \Delta \cdot
    \max\!\bigl(R(T_C, T_B; X_{T_C}) - K,\, 0\bigr).
  \label{eq:terminal}
\end{equation}
The model price at $t = 0$ is
\begin{equation}
  V_0 = \E^{\Q}\!\left[\exp\!\Bigl(-\int_0^{T_C} r(s)\dd s\Bigr)
         \phi(X_{T_C})\right].
  \label{eq:price}
\end{equation}

\subsection{The backward SDE}

By the Feynman--Kac theorem, $V_t \equiv \widehat{V}(t, X_t)$ satisfies the
forward-backward SDE (FBSDE): the forward component is \eqref{eq:X} and the
backward component is
\begin{equation}
  \widehat{V}_t = \phi(X_{T_C})
    + \int_t^{T_C} r(s, X_s)\, \widehat{V}_s \dd s
    - \int_t^{T_C} \widehat{Z}_s^\top \dW_s,
  \label{eq:bsde}
\end{equation}
where $r(t, X_t) = f_0 + \sum_{i=1}^4 X_i(t)$ and the control process
\begin{equation}
  \widehat{Z}_t = \nabla_X \widehat{V}(t, X_t)^\top \sigma(t) \in \R^3
  \label{eq:Z}
\end{equation}
is the $3$-dimensional martingale integrand, one component per Brownian
motion.

Several structural properties of \eqref{eq:bsde} are decisive for the
numerical treatment.

\begin{remark}[Linear, non-recursive driver]
The driver $h(t, x, y, z) = -r(t, x)\, y$ is linear in $y$ and independent
of $z$. In the terminology of \cite{GPR2022}, \S3.2, this is the
\emph{non-recursive} case: the backward component does not feed back into the
forward dynamics, so Algorithm~1 of \cite{GPR2022} applies directly without
the recursive nesting required for FVA-type adjustments.
\end{remark}

\begin{remark}[Lipschitz driver]
Since $r(t, X_t)$ is bounded on bounded sets, $h$ is uniformly Lipschitz in
$y$ on any compact region. The a priori error bounds of Han \& Long
\cite{HL2020}, Theorem~1', therefore apply and ensure that
\begin{equation}
  \sup_{t \in [0,T_C]} \E\bigl|\widehat{V}_t - \widehat{V}^{\,\xi,\rho}_t\bigr|^2
  + \int_0^{T_C} \E\bigl|\widehat{Z}_t - Z^{\rho}_t\bigr|^2 \dd t
  \;\leq\; C\left(\Delta t
    + \E\bigl|\phi(X_{T_C}) - \widehat{V}^{\,\xi,\rho}_{T_C}\bigr|^2\right),
  \label{eq:error_bound}
\end{equation}
where $C$ is independent of $\Delta t$ and the network parameters
$(\xi, \rho)$.
\end{remark}

% ─────────────────────────────────────────────────────────────────────────────
\section{Deep BSDE Solver}
\label{sec:solver}

\subsection{Time discretisation}

Let $0 = t_0 < t_1 < \cdots < t_N = T_C$ be a uniform grid with step
$\Delta t = T_C / N$. The Euler--Maruyama scheme for the forward component is
\begin{equation}
  X_{n+1} = X_n + \mu(t_n, X_n)\Delta t + \sigma(t_n)\,\Delta W_n,
  \qquad X_0 = 0,
  \label{eq:euler_forward}
\end{equation}
where $\Delta W_n = W_{t_{n+1}} - W_{t_n} \sim \mathcal{N}(0, \Delta t\, I_3)$.
The backward Euler scheme is
\begin{equation}
  \widehat{V}^{\,\xi,\rho}_{n+1} = \widehat{V}^{\,\xi,\rho}_n
    + r(t_n, X_n)\,\widehat{V}^{\,\xi,\rho}_n \Delta t
    + (Z^\rho_n)^\top \Delta W_n,
  \qquad \widehat{V}^{\,\xi,\rho}_0 = \xi.
  \label{eq:euler_backward}
\end{equation}
The key observation is that the $+r\,V\,\Delta t$ sign in \eqref{eq:euler_backward}
follows from $-h\,\Delta t = +r\,V\,\Delta t$ (the BSDE driver is $h = -rV$).

\subsection{Neural network parametrisation}

Following \cite{HJE2018}, the control $Z^\rho_n \approx \widehat{Z}_{t_n}$ is
parametrised by a feedforward network. We use a \emph{shared-weight} architecture:
a single MLP $\varphi^\rho: \R^5 \to \R^3$ with input $(t_n, X_n) \in \R^5$
is evaluated at each time step, rather than training a distinct network per step.
This is more parsimonious than the original Han et al.\ formulation and is
appropriate here because $\widehat{Z}(t, x)$ is a smooth function of $t$ under
the Cheyette dynamics, so a shared architecture captures the time dependence
without per-step expressivity.

The network has three hidden layers of width $32$ (i.e.\ $d + 28$, where
$d = 4$ is the state dimension), batch normalisation between layers, and
$\tanh$ activations. The output layer maps to $\R^3$. The full parameter
count is $\mathcal{O}(5 \cdot 32 + 32^2 \cdot 2 + 32 \cdot 3) \approx 2\,300$
weights.

The choice of $\tanh$ over ReLU is deliberate: the caplet price function is
smooth away from the exercise boundary, and piecewise-linear activations
impose artificial kinks that slow convergence toward the tight target.
Batch normalisation is critical for stable training, as noted in \cite{HJE2018}.

\subsection{Optimisation}

The parameters $(\xi, \rho)$ — comprising the scalar initial value $\xi \in \R$
and the network weights $\rho$ — are optimised jointly by minimising the
empirical terminal loss over $L$ Monte Carlo paths:
\begin{equation}
  \mathcal{L}(\xi, \rho) = \frac{1}{L}\sum_{\ell=1}^{L}
    \bigl(\phi(X^{(\ell)}_{T_C}) - \widehat{V}^{\,\xi,\rho,(\ell)}_{N}\bigr)^2.
  \label{eq:loss}
\end{equation}
We use Adam with a three-phase learning rate schedule:
\begin{equation}
  \eta =
  \begin{cases}
    5 \times 10^{-3} & \text{iterations } 1\text{--}5{,}000, \\
    5 \times 10^{-4} & \text{iterations } 5{,}001\text{--}10{,}000, \\
    5 \times 10^{-5} & \text{iterations } 10{,}001\text{--}20{,}000.
  \end{cases}
  \label{eq:lr}
\end{equation}
The final phase at $\eta = 5 \times 10^{-5}$ is the critical addition over
the two-phase schedule of \cite{HJE2018}; it provides the fine-grained descent
needed for the last decimal place of accuracy.

Hyperparameters are summarised in Table~\ref{tab:hypers}.

\begin{table}[h]
\centering
\begin{tabular}{lll}
\toprule
Parameter & Value & Rationale \\
\midrule
Time steps $N$ & $50 \cdot T_C$ & Euler strong error $\sim \Delta t^{1/2}$; 50 steps/year budgets $\sim$1\% per-step error \\
Training paths $L$ & $4{,}096$ & $4\times$ the Han et al.\ baseline for lower MC variance \\
Batch size $B$ & $256$ & Larger batches reduce gradient noise \\
Iterations $\mathcal{I}$ & $20{,}000$ & Three-phase LR; diminishing returns past $10^4$ \\
Hidden layers & $3$ & One additional layer beyond \cite{HJE2018} for tighter target \\
Width & $32$ & Balanced capacity vs.\ over-parameterisation \\
Activation & $\tanh$ & Smooth pricing function; ReLU too coarse \\
Precision & \texttt{float64} & Tolerance is $\sim$4 significant figures below price level \\
\bottomrule
\end{tabular}
\caption{Solver hyperparameters for the $0.1$ bps vega-adjusted target.}
\label{tab:hypers}
\end{table}

\subsection{Variance reduction}

\paragraph{Antithetic sampling.}
For each batch, $L/2$ independent Gaussian increments $\{Z_n^{(\ell)}\}_{\ell=1}^{L/2}$
are drawn and paired with their negatives $\{-Z_n^{(\ell)}\}$. This halves
the estimator variance with no additional network evaluations.

\paragraph{Warm-start initialisation.}
Remark~4.1 of \cite{GPR2022} establishes that the error in $V_0$ is
bottlenecked by $\xi$, decoupled from $\rho$: the network controls the
hedging accuracy but $\xi$ alone sets the initial price estimate. It is
therefore natural to initialise $\xi$ near the true price.

A tempting but \emph{invalid} warm-start is to use the analytical Cheyette
integrated variance $v^2(0, T_C, T_B)$ in a Black formula, since the
Cheyette caplet formula is precisely Black with that $v$. This would
initialise $\xi$ at the analytical price to machine precision, making the
BSDE solver redundant.

The correct approach is an \emph{honest} warm-start: the standard textbook
Black caplet on LIBOR with a flat $1\%$ vol prior,
\begin{equation}
  \xi_0 = e^{-f_0 T_B} \Delta
    \bigl[F \cdot \Phi(d_1) - K \cdot \Phi(d_2)\bigr], \qquad
  d_{1,2} = \frac{\ln(F/K) \pm \tfrac{1}{2}\sigma_0^2 T_C}{\sigma_0\sqrt{T_C}},
  \label{eq:warmstart}
\end{equation}
with $\sigma_0 = 0.01$ and $F = \Delta^{-1}(e^{f_0\Delta} - 1)$ the forward
LIBOR rate. For the 2Y ATM caplet under calibrated Cheyette parameters this
gives $\xi_0 \approx 3.1 \times 10^{-4}$, roughly $89\%$ below the true
price of $2.79 \times 10^{-3}$. The solver must close this gap.

% ─────────────────────────────────────────────────────────────────────────────
\section{Accuracy Target and Error Metric}
\label{sec:accuracy}

\subsection{Vega-adjusted error}

Raw price errors are not a meaningful accuracy metric across expiries and
strikes because the price level varies by an order of magnitude across the
grid. The natural unit is the implied-vol error, obtained by dividing the
price error by the caplet's Black vega:
\begin{equation}
  \varepsilon_{\nu} = \frac{\widehat{V}_0 - V_0^{\mathrm{analytical}}}{\nu},
  \qquad
  \nu = B(0, T_C)\,\phi(d_1)\sqrt{T_C},
  \label{eq:vegaerr}
\end{equation}
where $\phi$ is the standard normal density and $d_1$ is evaluated at the
analytical implied vol. The target is $|\varepsilon_\nu| < 10^{-5}$
(= $0.1$ basis points).

\subsection{Budget across the grid}

Table~\ref{tab:budget} translates the $0.1$ bps target into absolute price
tolerances for the test grid. The raw tolerance of $\mathcal{O}(10^{-6})$
across all expiries sits approximately four decimal places below the price
level, which is roughly \emph{two orders of magnitude tighter} than the
benchmark errors reported in \cite{HJE2018} for comparable European payoffs.

\begin{table}[h]
\centering
\begin{tabular}{cccccc}
\toprule
$T_C$ (Y) & Price & IV & Vega & $|\varepsilon_\nu| < 0.1$ bps $\Rightarrow$ & $|\Delta V| <$ \\
\midrule
1 & $2.07 \times 10^{-3}$ & $0.72\%$ & $3.79 \times 10^{-1}$ & & $3.79 \times 10^{-6}$ \\
2 & $2.79 \times 10^{-3}$ & $1.04\%$ & $5.10 \times 10^{-1}$ & & $5.10 \times 10^{-6}$ \\
3 & $3.29 \times 10^{-3}$ & $1.30\%$ & $5.95 \times 10^{-1}$ & & $5.95 \times 10^{-6}$ \\
5 & $3.99 \times 10^{-3}$ & $1.76\%$ & $6.94 \times 10^{-1}$ & & $6.94 \times 10^{-6}$ \\
\bottomrule
\end{tabular}
\caption{ATM caplet reference values (6M tenor, flat curve $f_0 = 5\%$,
calibrated Cheyette parameters) and absolute price tolerances implied by the
$0.1$ bps vega-adjusted target.}
\label{tab:budget}
\end{table}

% ─────────────────────────────────────────────────────────────────────────────
\section{Analytical Benchmark}
\label{sec:analytical}

The closed-form caplet price in the Cheyette model \cite{BCK2012} is
\begin{equation}
  V_0 = B(0, T_C)\Phi(d_1) - B(0, T_B)(1 + K\Delta)\Phi(d_2),
  \label{eq:analytical}
\end{equation}
with
\begin{equation}
  d_{1,2} = \frac{\ln\!\bigl[B(0,T_C) / B(0,T_B)(1+K\Delta)\bigr]
                  \pm \tfrac{1}{2}\,v^2}{v},
\end{equation}
where $v^2 \equiv v^2(0, T_C, T_B)$ is the model-implied integrated variance
computed from the three-factor volatility structure via explicit quadratures
over the $G$ functions (equations U1--U3 of \cite{BCK2012}).

Two consequences follow. First, the closed-form price provides an exact
benchmark for the deep BSDE output, enabling the tight error metric of
\eqref{eq:vegaerr}. Second — and critically for the initialisation strategy
— the formula \eqref{eq:analytical} is \emph{identical} to the Black formula
with $v$ set to the Cheyette integrated vol. Any warm-start that uses
$v^2(0, T_C, T_B)$ to initialise $\xi$ therefore recovers the analytical
price directly, bypassing the solver entirely.

% ─────────────────────────────────────────────────────────────────────────────
\section{Error Analysis Grid}
\label{sec:grid}

The solver is tested on a $4 \times 3$ grid of expiry--moneyness pairs:
$T_C \in \{1, 2, 3, 5\}$ years with 6M tenor, and moneyness
$m = K / K_{\mathrm{ATM}} \in \{0.7, 1.0, 1.3\}$. The deep ITM and OTM cells
probe the failure mode that most concerns us: the network must learn near-zero
output (deep OTM) or nearly intrinsic output (deep ITM) through smooth
interpolation away from the payoff kink at $R = K$.

For each cell, the solver is trained independently from the honest warm-start
\eqref{eq:warmstart}. The vega-adjusted error \eqref{eq:vegaerr} is reported
against the closed-form benchmark \eqref{eq:analytical}. The pass criterion
is $|\varepsilon_\nu| < 10^{-5}$ uniformly across the grid.

% ─────────────────────────────────────────────────────────────────────────────
\section{Extensions and Limitations}
\label{sec:extensions}

\paragraph{Control variate.}
If the $0.1$ bps target is not achieved uniformly, a natural improvement is
to subtract the bond price $B(T_C, T_B; X_{T_C})$ from the terminal payoff
and add back its known analytical expectation. The bond price has a closed
form under the Cheyette model, so the control variate is exact. The residual
payoff $\phi - B$ is zero at-the-money by construction, reducing the
effective variance of \eqref{eq:loss} by a factor of $2$--$4$.

\paragraph{Quasi-Monte Carlo.}
Replacing the pseudo-random Brownian increments with a scrambled Sobol
sequence in $\R^{3N}$ (one dimension per Brownian per timestep) reduces the
MC variance at rate $\mathcal{O}((\log L)^{3N}/L)$ vs $\mathcal{O}(L^{-1/2})$.
For $N = 100$ and $d = 3$ the effective dimension is $300$, at which point QMC
gains are modest unless effective dimension reduction (e.g.\ Brownian bridge
construction) is applied.

\paragraph{Bermudan swaptions and xVA.}
The caplet FBSDE studied here is the simplest non-trivial case. The same
architecture extends to Bermudan swaptions by adding an early-exercise
constraint in the backward propagation (cf.\ \cite{GPR2022} for the
primal-dual approach). For xVA computations, the clean caplet price
trajectory produced by the trained solver feeds directly into the outer
CVA/DVA integral via Algorithm~2 of \cite{GPR2022}, or into the recursive
FVA BSDE via Algorithm~3.

\paragraph{Limitations.}
The Euler--Maruyama scheme introduces a strong error of order
$\mathcal{O}(\Delta t^{1/2})$. With $N = 50\, T_C$ steps, this error is
approximately $1\%$ of the state standard deviation per step, subdominant to
the target accuracy. Increasing $N$ beyond $100$ per year gives diminishing
returns as the MC variance in the terminal loss begins to dominate. For longer
expiries ($T_C > 5$Y), a Milstein scheme or exact simulation of the Gaussian
marginals would be preferable.

% ─────────────────────────────────────────────────────────────────────────────
\section*{References}

\begin{thebibliography}{9}

\bibitem{BCK2012}
Beyna, I., Chiarella, C., \& Kang, B.\ (2012).
Pricing interest rate derivatives in a multifactor HJM model with time
dependent volatility.
\textit{SSRN Working Paper} 2162748.

\bibitem{HJE2018}
Han, J., Jentzen, A., \& E, W.\ (2018).
Solving high-dimensional partial differential equations using deep learning.
\textit{Proceedings of the National Academy of Sciences},
115(34), 8505--8510.

\bibitem{HL2020}
Han, J., \& Long, J.\ (2020).
Convergence of the deep BSDE method for coupled FBSDEs.
\textit{Probability, Uncertainty and Quantitative Risk}, 5(1), 1--33.

\bibitem{GPR2022}
Gnoatto, A., Picarelli, A., \& Reisinger, C.\ (2022).
Deep xVA solver --- a neural network based counterparty credit risk management
framework.
\textit{arXiv:2005.02633v4}.

\end{thebibliography}

\end{document}
